% Generated draft: requires mathematical and source review.
\documentclass[11pt]{article}
\usepackage[margin=25mm]{geometry}
\usepackage{fontspec}
\setmainfont{Latin Modern Roman}
\usepackage{amsmath,amssymb,mathtools,unicode-math}
\setmathfont{Latin Modern Math}
\usepackage{enumitem,needspace,xurl,hyperref}
\setlength{\parindent}{0pt}
\setlength{\parskip}{5pt}
\setlength{\emergencystretch}{3em}
\newcommand{\N}{\mathbb{N}}
\newcommand{\Z}{\mathbb{Z}}
\newcommand{\Q}{\mathbb{Q}}
\newcommand{\R}{\mathbb{R}}
\newcommand{\C}{\mathbb{C}}
\newcommand{\F}{\mathbb{F}}
\newcommand{\A}{\mathbb{A}}
\newcommand{\PP}{\mathbb{P}}
\newcommand{\E}{\mathbb{E}}
\newcommand{\eps}{\varepsilon}
\newcommand{\dd}{\,\mathrm{d}}
\begin{document}
\section*{UnsolvedMath ID: 30001495}
\textbf{Source problem code:} catalogue entry
\subsection{Definitions and mathematical statement}
Let \(X\hookrightarrow\PP^N\) be a smooth \(n\)-dimensional variety, with \(N=n+c\). The ground field and its characteristic are not specified in the supplied record. Write
\[
\mathcal J_X=\ker\bigl(\mathcal O_{\PP^N}\longrightarrow\mathcal O_X\bigr),\qquad
\mathcal J_X(q)=\mathcal J_X\otimes\mathcal O_{\PP^N}(q).
\]
Here \(H^p(\mathcal J_X(q))\) means \(H^p(\PP^N,\mathcal J_X(q))\). A cohomological degree is a nonnegative integer, and a twist is an integer; the record supplies no further ranges for \(p,q\), or \(c\).

For a general projection of \(X\), let \(X_k\) denote the locus called the ``multiplicity-at-least-\(k\) locus'' in the source record, and write \(\widetilde X_k\) for its normalization, for positive integers \(k\). This notation remains provisional: the projection target and center, the meaning of multiplicity, the scheme structure of \(X_k\), and the normalization convention are not specified. In particular, no particular projection or definition of this locus is selected here.

The supplied conjectural package asks whether the following assertions hold for every \(X\) as above and all relevant indices:
\begin{enumerate}
\item For every nonnegative integer \(p\) and integer \(q\) in the intended, unspecified index ranges,
\[
p+(q+1)(c-1)\leq n\quad\text{and}\quad p<N-1
\quad\Longrightarrow\quad
H^p(\PP^N,\mathcal J_X(q))=0.
\]
\item For every positive integer \(k\),
\[
k(c-1)<n
\quad\Longrightarrow\quad
\widetilde X_k\text{ is smooth and connected}.
\]
\item For every positive integer \(k\),
\[
k(c-1)\leq n\quad\text{and}\quad\widetilde X_k\text{ is connected}
\quad\Longrightarrow\quad
H^1(\PP^N,\mathcal J_X(k-1))=0.
\]
\end{enumerate}
The strict inequality in assertion (2), the weak inequalities in assertions (1) and (3), and the bound \(p<N-1\) are retained exactly. These are source assertions requiring review; missing hypotheses and projection data prevent a fully determined formulation.
\subsection{Short English statement}
For a smooth projective variety, the supplied conjectural package relates ideal-sheaf cohomology vanishing to smoothness and connectedness of normalized multiple-point loci of a general projection. The ground field, some index restrictions, and the precise projection and multiple-point construction remain unspecified.
\subsection{Sources}
\begin{itemize}
\item Supplied local UnsolvedMath registry record, permanent ID \(30001495\), exact source problem code \texttt{catalogue entry}, title ``Projection-Locus Smoothness and Cohomology Vanishing,'' in the Oberwolfach Reports Open Problems collection. The supplied proposed year is 2010; no individual proposer is identified.
\item Supplied report citation: \emph{Classical Algebraic Geometry} (2010), Oberwolfach Reports, DOI 10.4171/owr/2010/27, pp.\ 1573--1623. This is also the supplied literature reference ``original OWR Report 27/2010'': \url{https://doi.org/10.4171/owr/2010/27}. The background describes these reports as records of mathematical workshops and their open-problem discussions.
\item Supplied literature reference, ``Gruson--Peskine general projection theorem'': \url{https://arxiv.org/abs/1010.2399}.
\item Supplied literature reference, ``Ran on generic projection and filling secants'': \url{https://doi.org/10.1215/00127094-2881483}.
\item Literature assessment dated exactly \texttt{2026-08-21}: the supplied status is ``partially solved.'' The supplied assessment attributes the principal smoothness and expected-dimension mechanism to the Gruson--Peskine general projection theorem, mentions later extensions, and states that the full three-part package, including all connectedness and vanishing assertions, is not established as one theorem. This is supplied dated curation and literature triage, not an independently verified assessment; its status and cited claims require independent verification.
\item The registry labels the statement \texttt{corrected\_verified} and reports that it was checked against the official Oberwolfach report. That verification claim is supplied provenance, not verification performed for this draft.
\end{itemize}
\end{document}
